% !PRINTED 348
\weberpage{348}

%<<CONT
\gesperrt{ist, bestimmt.} Es muß also vier ganze rationale Zahlen
$\alpha$, $\beta$, $\gamma$, $\delta$ geben, die der Bedingung
\weq{2}{\begin{aligned}\eta\, \mathfrak{a}' &= \mathfrak{a}\,(\alpha - \gamma\,\omega),\\ \eta\, \mathfrak{a}'\,\omega' &= \mathfrak{a}\,(-\beta + \delta\,\omega)\end{aligned}}
genügen. Bezeichnen wir für den Augenblick mit $\eta_1$, $\omega_1$, $\omega_1'$ die
zu $\eta$, $\omega$, $\omega'$ konjugierten Zahlen, so ist hiernach
\[
\mathfrak{a}'^2\, \eta\, \eta_1
\begin{vmatrix} 1, & \omega' \\ 1, & \omega_1' \end{vmatrix}
= \mathfrak{a}^2
\begin{vmatrix} \alpha, & -\gamma \\ -\beta, & \delta \end{vmatrix}
\begin{vmatrix} 1, & \omega \\ 1, & \omega_1 \end{vmatrix},
\]
also:
\[
\mathfrak{a}'^2\, \eta\, \eta_1\,(\omega' - \omega_1') = \mathfrak{a}^2\,(\alpha\delta - \beta\gamma)\,(\omega - \omega_1),
\]
und da nach \S\,91, (12):
\[
\mathfrak{a}\,(\omega - \omega_1) = \mathfrak{a}'\,(\omega' - \omega_1') = -\sqrt{\Delta}
\]
ist, so folgt hieraus, wenn
\[
\varepsilon = \alpha\delta - \beta\gamma
\]
gesetzt wird,
\weq{3}{\mathfrak{a}\,\varepsilon = \mathfrak{a}'\,\eta\,\eta_1;}
wenn also, wie wir schon angenommen haben, $\mathfrak{a}$ und $\mathfrak{a}'$ positiv
sind und (nach dem verschärften Äquivalenzbegriff) $\eta\eta_1$ positiv
ist, so ist auch $\varepsilon$ positiv.

Vertauschen wir aber $\mathfrak{a}$ mit $\mathfrak{a}'$, so geht $\eta$ in $1 : \eta$ über und
$\varepsilon$ mag in $\varepsilon'$ übergehen. Demnach folgt auch
\weq{4}{\mathfrak{a}'\,\varepsilon' = \frac{\mathfrak{a}}{\eta\,\eta_1},}
und folglich $\varepsilon \varepsilon' = 1$. Da aber $\varepsilon$ und $\varepsilon'$ positive ganze rationale
Zahlen sind, so folgt hieraus $\varepsilon = \varepsilon' = 1$. Also nach (2)
\weq{5}{\omega' = \frac{-\beta - \delta\,\omega}{\alpha - \gamma\,\omega}, \qquad \omega = \frac{\alpha\,\omega' + \beta}{\gamma\,\omega' + \delta},}
\[
\alpha\delta - \beta\gamma = 1.
\]

Nun sind $\omega$ und $\omega'$ die Wurzeln der beiden Formen
\[
(a,\, b,\, c), \quad (a',\, b',\, c') \qquad [\S\,91,\ (13)],
\]
und durch (5) ist die Äquivalenz dieser Formen ausgedrückt.
Wir haben daher, übereinstimmend mit der allgemeinen Theorie
(\S\,88) den Satz:

5. \gesperrt{Die zu äquivalenten Idealen gehörigen quadratischen
Formen sind äquivalent.}

Man erkennt hieraus die Bedeutung des verschärften Äquivalenzbegriffes
: Wollte man nämlich auch negative Werte von
$N(\eta)$ zulassen, so könnten auch $\varepsilon = -1$ und die entsprechenden
Formen \gesperrt{uneigentlich} äquivalent sein.
